\documentclass[addpoints]{exam}
\usepackage{amsmath}
\usepackage{amssymb}
\usepackage{color}

% \printanswers

\newcommand{\event}{Exponential Functions}
\newcommand{\tournament}{}
\def\type{0}

\setlength\fillinlinelength{\textwidth}
\CorrectChoiceEmphasis{\color{red}\bfseries}
\SolutionEmphasis{\color{red}\bfseries}

\setlength\answerclearance{2pt}
\pagestyle{headandfoot}
\runningheadrule
\runningheader{\event}{\tournament}{\ifnum\type=1 Class Set Test \else Full Name:\hspace{1cm} \fi}
\runningfootrule
\runningfooter{}{Page \thepage\ of \numpages}{}

\begin{document}
\begin{center}
    {\huge Grade 11 Foundations for College Math \\ \event
    \ifprintanswers \space Key
    \else \space \ifnum\type<2 Test \fi \ifnum\type=1 \textbf{(CLASS SET)} \fi
          \ifnum\type=2 Answer Sheet \fi
    \fi \\ \vspace{3mm}\tournament\vspace{1mm}}
\end{center}

\vspace{.25\textheight}

\begin{center}
    First Name: \ifprintanswers
    \underline{\hspace{3.5cm}}\textbf{KEY}\underline{\hspace{5.5cm}}\\\ \vspace{5mm}
    \else \underline{\hspace{10cm}}\\ \vspace{5mm} \fi
    Last Name: \underline{\hspace{10cm}}\\ \vspace{5mm}
\end{center}

\noindent\textbf{\underline{Directions:}}
\begin{itemize}
    \ifnum\type=1 \item \textbf{This test is a class set.} Do not write on this paper. \fi
    \item Show all work for full marks. Round dollar amounts to the nearest cent.
    \item The test is designed to be completed in 75 minutes.
\end{itemize}
\begin{center}
\textit{For grading use only}\\
\multirowgradetable{1}[pages]
\end{center}

\newpage
\begin{questions}

\section*{Part A --- Multiple Choice (15 marks)}
\vspace{5mm}

% Q1
\question[1]
Evaluate $2^{-3}$.
\begin{choices}
    \choice $-8$
    \choice $-6$
    \correctchoice $\dfrac{1}{8}$
    \choice $\dfrac{1}{6}$
\end{choices}

% Q2
\question[1]
Which expression equals $27^{2/3}$?
\begin{choices}
    \choice $9$
    \correctchoice $9$
    \choice $3$
    \choice $81$
\end{choices}

% Q3 — rewrite to avoid duplicate answer issue
\question[1]
Simplify $27^{1/3}$.
\begin{choices}
    \choice $9$
    \correctchoice $3$
    \choice $27$
    \choice $81$
\end{choices}

% Q4
\question[1]
The function $y = 3(2)^x$ represents:
\begin{choices}
    \choice Exponential decay, because the base is 2.
    \correctchoice Exponential growth, because the base $2 > 1$.
    \choice A linear function.
    \choice Exponential decay, because the coefficient is 3.
\end{choices}

% Q5
\question[1]
Which of the following is an exponential decay function?
\begin{choices}
    \choice $y = 4(1.1)^x$
    \choice $y = 2(3)^x$
    \correctchoice $y = 5(0.6)^x$
    \choice $y = (1)^x$
\end{choices}

% Q6
\question[1]
For the function $y = a(b)^x$, the $y$-intercept is:
\begin{choices}
    \choice $b$
    \choice $0$
    \correctchoice $a$
    \choice $ab$
\end{choices}

% Q7
\question[1]
What is the horizontal asymptote of $y = 4(2)^x$?
\begin{choices}
    \choice $y = 4$
    \correctchoice $y = 0$
    \choice $y = 2$
    \choice $y = 1$
\end{choices}

% Q8
\question[1]
A city has 50{,}000 people and grows at 3\% per year. After 1 year the population is:
\begin{choices}
    \choice 51{,}500
    \correctchoice 51{,}500
    \choice 53{,}000
    \choice 50{,}300
\end{choices}

% Q8 — fix duplicate
\question[1]
A city of 50{,}000 grows at 3\% per year. Which equation models the population $P$ after $t$ years?
\begin{choices}
    \choice $P = 50000 + 0.03t$
    \correctchoice $P = 50000(1.03)^t$
    \choice $P = 50000(0.03)^t$
    \choice $P = 50000(1.3)^t$
\end{choices}

% Q9
\question[1]
A substance has a half-life of 10 years. Starting with 80 g, how much remains after 20 years?
\begin{choices}
    \choice 40 g
    \correctchoice 20 g
    \choice 10 g
    \choice 5 g
\end{choices}

% Q10
\question[1]
Solve for $x$: $2^x = 32$.
\begin{choices}
    \choice $x = 4$
    \correctchoice $x = 5$
    \choice $x = 6$
    \choice $x = 16$
\end{choices}

% Q11
\question[1]
Solve for $x$: $3^{2x} = 81$.
\begin{choices}
    \choice $x = 4$
    \choice $x = 3$
    \correctchoice $x = 2$
    \choice $x = 1$
\end{choices}

% Q12
\question[1]
The graph of $y = 2^x$ is shifted 3 units to the right. The new equation is:
\begin{choices}
    \choice $y = 2^x + 3$
    \choice $y = 2^{x+3}$
    \correctchoice $y = 2^{x-3}$
    \choice $y = 2^x - 3$
\end{choices}

% Q13
\question[1]
What is the effect of changing $a$ in $y = a(b)^x$ from 1 to 4 (with $b > 1$)?
\begin{choices}
    \choice The asymptote moves to $y = 4$.
    \correctchoice The $y$-intercept shifts from 1 to 4 and the graph is vertically stretched.
    \choice The base changes.
    \choice The function becomes a decay function.
\end{choices}

% Q14
\question[1]
A bacteria culture doubles every 3 hours. Which expression gives the number after $n$ doubling periods starting from 500?
\begin{choices}
    \choice $500 + 2n$
    \correctchoice $500 \cdot 2^n$
    \choice $500 \cdot 3^n$
    \choice $2 \cdot 500^n$
\end{choices}

% Q15
\question[1]
An exponential function passes through $(0, 6)$ and $(1, 18)$. What is the equation?
\begin{choices}
    \choice $y = 18(6)^x$
    \correctchoice $y = 6(3)^x$
    \choice $y = 3(6)^x$
    \choice $y = 6(18)^x$
\end{choices}

\section*{Part B --- Short Answer (33 marks)}

% Q16
\question[6]
\textbf{Evaluating Exponential Expressions.} Evaluate each expression. Show all steps.

\begin{parts}
    \part[2] $4^{-2} + 8^{0}$
    \begin{solution}
        $4^{-2} = \dfrac{1}{4^2} = \dfrac{1}{16}$\\
        $8^{0} = 1$\\
        $4^{-2} + 8^{0} = \dfrac{1}{16} + 1 = \dfrac{17}{16}$ or $1.0625$
    \end{solution}

    \part[2] $16^{3/4}$
    \begin{solution}
        $16^{3/4} = \left(16^{1/4}\right)^3 = 2^3 = 8$
    \end{solution}

    \part[2] Simplify: $\dfrac{3^5 \cdot 3^{-2}}{3^2}$
    \begin{solution}
        $\dfrac{3^5 \cdot 3^{-2}}{3^2} = \dfrac{3^{5+(-2)}}{3^2} = \dfrac{3^3}{3^2} = 3^{3-2} = 3^1 = 3$
    \end{solution}
\end{parts}

% Q17
\question[6]
\textbf{Identifying Growth and Decay.} For each function below, state whether it represents growth or decay, identify the values of $a$ and $b$, and state the $y$-intercept and horizontal asymptote.

\begin{parts}
    \part[3] $y = 7(1.4)^x$
    \begin{solution}
        \textbf{Growth} (base $b = 1.4 > 1$).\\
        $a = 7$, $b = 1.4$.\\
        $y$-intercept: $(0, 7)$ (substitute $x = 0$: $y = 7(1.4)^0 = 7$).\\
        Horizontal asymptote: $y = 0$.
    \end{solution}

    \part[3] $y = 12(0.75)^x$
    \begin{solution}
        \textbf{Decay} (base $b = 0.75$, and $0 < 0.75 < 1$).\\
        $a = 12$, $b = 0.75$.\\
        $y$-intercept: $(0, 12)$.\\
        Horizontal asymptote: $y = 0$.
    \end{solution}
\end{parts}

% Q18
\question[6]
\textbf{Population Growth.} A city has a current population of 80{,}000 and is growing at a rate of 2.5\% per year.

\begin{parts}
    \part[2] Write an equation to model the population $P$ after $t$ years.
    \begin{solution}
        $P = 80000(1.025)^t$
    \end{solution}

    \part[2] Find the population after 10 years. Round to the nearest person.
    \begin{solution}
        $P = 80000(1.025)^{10}$\\
        $(1.025)^{10} \approx 1.28008$\\
        $P \approx 80000 \times 1.28008 \approx 102{,}406$ people
    \end{solution}

    \part[2] Estimate when the population will reach 120{,}000 by testing values of $t$ in your equation. Show at least two test values.
    \begin{solution}
        We need $80000(1.025)^t = 120000$, i.e.\ $(1.025)^t = 1.5$.\\
        Test $t = 16$: $(1.025)^{16} \approx 1.4845$ (too low)\\
        Test $t = 17$: $(1.025)^{17} \approx 1.5216$ (too high)\\
        The population reaches 120{,}000 between year 16 and year 17, so approximately \textbf{17 years} from now.
    \end{solution}
\end{parts}

% Q19
\question[7]
\textbf{Radioactive Decay.} A radioactive substance has a half-life of 12 years. You start with 200 g.

\begin{parts}
    \part[2] Write an equation for the mass $M$ (in grams) remaining after $t$ years.
    \begin{solution}
        Each 12 years, the amount is multiplied by $\tfrac{1}{2}$. The number of half-lives in $t$ years is $\dfrac{t}{12}$.
        \[M = 200\!\left(\tfrac{1}{2}\right)^{t/12}\]
    \end{solution}

    \part[2] Find the amount remaining after 36 years.
    \begin{solution}
        $t = 36$: number of half-lives $= \dfrac{36}{12} = 3$.
        \[M = 200\!\left(\tfrac{1}{2}\right)^{3} = 200 \times \tfrac{1}{8} = 25 \text{ g}\]
    \end{solution}

    \part[3] Sketch the graph of $M$ vs.\ $t$ for $0 \le t \le 48$ years. Label the points at $t = 0, 12, 24, 36, 48$ with their coordinates, and draw the asymptote.
    \begin{solution}
        Key points:
        \begin{itemize}
            \item $(0,\; 200)$
            \item $(12,\; 100)$
            \item $(24,\; 50)$
            \item $(36,\; 25)$
            \item $(48,\; 12.5)$
        \end{itemize}
        The curve decreases steeply at first then levels off, approaching the horizontal asymptote $M = 0$ (but never reaching it). Students should draw a smooth decreasing curve through all five labeled points.
    \end{solution}
\end{parts}

% Q20
\question[8]
\textbf{Comparing Two Investment Models.} Two investment options are modelled as follows:
\begin{itemize}
    \item \textbf{Option A:} $V_A = 1000(1.08)^t$ (8\% annual growth)
    \item \textbf{Option B:} $V_B = 1500(1.05)^t$ (5\% annual growth, higher start)
\end{itemize}
where $t$ is time in years.

\begin{parts}
    \part[2] What is the initial value (at $t = 0$) of each option?
    \begin{solution}
        Option A: $V_A(0) = 1000(1.08)^0 = \$1{,}000$\\
        Option B: $V_B(0) = 1500(1.05)^0 = \$1{,}500$
    \end{solution}

    \part[3] Find the value of each option at $t = 5$, $t = 10$, and $t = 20$ years. Record results in a table.
    \begin{solution}
        \begin{tabular}{ccc}
        \hline
        $t$ & Option A & Option B\\
        \hline
        5  & $1000(1.08)^5 \approx \$1{,}469.33$  & $1500(1.05)^5 \approx \$1{,}914.42$ \\
        10 & $1000(1.08)^{10} \approx \$2{,}158.92$ & $1500(1.05)^{10} \approx \$2{,}443.74$ \\
        20 & $1000(1.08)^{20} \approx \$4{,}660.96$ & $1500(1.05)^{20} \approx \$3{,}979.95$ \\
        \hline
        \end{tabular}
    \end{solution}

    \part[3] Between which two years do the values become equal? Estimate by testing integer values of $t$. What does this tell us about which option is better?
    \begin{solution}
        We need $1000(1.08)^t = 1500(1.05)^t$.

        Test $t = 15$: $A \approx 1000(1.08)^{15} \approx 3172$, \; $B \approx 1500(1.05)^{15} \approx 3118$ $\Rightarrow A > B$.\\
        Test $t = 14$: $A \approx 1000(1.08)^{14} \approx 2937$, \; $B \approx 1500(1.05)^{14} \approx 2969$ $\Rightarrow A < B$.\\

        The values are equal between $t = 14$ and $t = 15$ years.\\
        \textbf{Interpretation:} Option B is better for the first $\approx 14$ years because it starts higher. After about 14--15 years, Option A overtakes B because its higher growth rate (8\%) compounds more aggressively over time.
    \end{solution}
\end{parts}

\end{questions}
\end{document}
